%% kapa: exam=2 round=36 year=2025 period=3 subject=compensation_law subject_ko=감정평가 및 보상법규 seq=2 number=2 points=30
%% source: Q-Net 감정평가사 기출문제 (cst00309) artlSeq=5239325 file=raw/1008/5239325/2025년 제36회 감정평가사 2차 감정평가 및 보상법규.pdf
%% conversion: kordoc 4.24.51 -> tools/gen2.py (auto); review=proofread-gpt-sol; flags=none
\begin{longquestion}{【문제 2】}
국토교통부장관은 X도 Y시 소재의 공부상 지목이 대(垈)인 A토지를 표준지로 선정하여 2024. 1. 1. 기준 표준지공시지가를 1m² 당 550만원으로 2024. 1. 25. 결정ㆍ공시하였다. 이에 따라 X도 Y시장은 Y시 부동산가격공시위원회의 심의를 거쳐 2024. 5. 31. 공부상 지목이 답(沓)인 甲 소유의 토지(이하 ‘이 사건 토지’라 함)의 개별공시지가를 1m² 당 85만원으로 결정ㆍ공시하고, 2024. 6. 3. 甲에게 이를 통지하였다. 다음 물음에 답하시오. (30점) \par
\medskip\noindent 물음 1) 이 사건 토지는 ｢국토의 계획 및 이용에 관한 법률｣에 따른 도시계획시설인 공원 부지로 토지수용의 대상이 되었다. 2024. 11. 15. 이 사건 토지에 대하여 관할 토지수용위원회의 수용재결이 이루어졌다. 이에 甲은 2024. 12. 2. 수용보상금증액청구소송을 제기하면서 표준지공시지가결정의 위법을 주장하고 있다. 甲의 주장이 인용될 수 있는지를 검토하시오. (20점) \par
\medskip\noindent 물음 2) 甲은 2024. 6. 14. 이 사건 토지의 개별공시지가결정에 대하여 ｢부동산 가격공시에 관한 법률｣에 따라 이의신청을 하였으나 2024. 6. 28. 기각결정의 통지를 받았다. 甲은 2024. 7. 30. 개별공시지가결정에 대하여 ｢행정심판법｣에 따른 행정심판을 제기하였으나, 2024. 8. 30. 기각재결을 받았고, 재결서 정본은 2024. 9. 10. 甲에게 송달되었다. 甲은 2024. 9. 20. 개별공시지가결정에 대한 취소소송을 제기하였다. 甲이 제기한 취소소송의 제소기간 준수 여부를 검토하시오. (10점) \par
\begin{framedtext}
참조 조문 \par
<부동산 가격공시에 관한 법률> \par
제11조(개별공시지가에 대한 이의신청) ① 개별공시지가에 이의가 있는 자는 그 결정ㆍ공시일부터 30일 이내에 서면으로 시장ㆍ군수 또는 구청장에게 이의를 신청할 수 있다. \par
② 시장ㆍ군수 또는 구청장은 제1항에 따라 이의신청 기간이 만료된 날부터 30일 이내에 이의신청을 심사하여 그 결과를 신청인에게 서면으로 통지하여야 한다. 이 경우 시장ㆍ군수 또는 구청장은 이의신청의 내용이 타당하다고 인정될 때에는 제10조에 따라 해당 개별공시지가를 조정하여 다시 결정ㆍ공시하여야 한다. \par
③ 제1항 및 제2항에서 규정한 것 외에 이의신청 및 처리절차 등에 필요한 사항은 대통령령으로 정한다. \par
<공익사업을 위한 토지 등의 취득 및 보상에 관한 법률> \par
제85조(행정소송의 제기) ① 사업시행자, 토지소유자 또는 관계인은 제34조에 따른 재결에 불복할 때에는 재결서를 받은 날부터 90일 이내에, 이의신청을 거쳤을 때에는 이의신청에 대한 재결서를 받은 날부터 60일 이내에 각각 행정소송을 제기할 수 있다. 이 경우 사업시행자는 행정소송을 제기하기 전에 \par
제84조에 따라 늘어난 보상금을 공탁하여야 하며, 보상금을 받을 자는 공탁된 보상금을 소송이 종결될 때까지 수령할 수 없다. \par
② 제1항에 따라 제기하려는 행정소송이 보상금의 증감(增減)에 관한 소송인 경우 그 소송을 제기하는 자가 토지소유자 또는 관계인일 때에는 사업시행자를, 사업시행자일 때에는 토지소유자 또는 관계인을 각각 피고로 한다. \par
<행정소송법> \par
제20조(제소기간) ① 취소소송은 처분 등이 있음을 안 날부터 90일 이내에 제기하여야 한다. 다만, \par
제18조 제1항 단서에 규정한 경우와 그 밖에 행정심판청구를 할 수 있는 경우 또는 행정청이 행정심판청구를 할 수 있다고 잘못 알린 경우에 행정심판청구가 있은 때의 기간은 재결서의 정본을 송달받은 날부터 기산한다. \par
\end{framedtext}
\end{longquestion}
